\subsection{凯莱代数}

定义 1. \(\mathbf{A}\) 上的一个凯莱代数是指一个有序对 \((\mathbf{E}, s)\) ，其中 \(\mathbf{E}\) 是 \(\mathbf{A}\) 上的一个带单位元 \(e\) 的代数，而 \(s\) 是 \(\mathbf{E}\) 的一个反自同构，使得

\[
u + s (u) \in \mathrm{A} e \quad \text { and } \quad u. s (u) \in \mathrm{A} e
\]

对所有 \(u \in \mathbf{E}\) 成立。

s 称为凯莱代数 (E, s) 的共轭，而 \(s(u)\) 称为 u 的共轭元。条件 \(u + s(u) \in Ae\) 蕴含 u 与 \(s(u)\) 可交换。我们记

\begin{enumerate}
\def\labelenumi{(\arabic{enumi})}
\setcounter{enumi}{8}
\tightlist
\item
\end{enumerate}

\[
\mathrm{T} (u) = u + s (u)\tag{10}
\]

\[
\mathrm{N} (u) = u. s (u) = s (u). u
\]

而子代数 Ae 的这些元素分别称为 u 的凯莱迹与范数。

由二次代数 E 及其共轭 s（因 E 交换，故 s 是反自同构）（第 3 号）组成的有序对是一个凯莱代数。

设 \((\mathbf{E}, s)\) 是一个凯莱代数；由于 \(s(e) = e\) ，有 \(s(u + s(u)) = u + s(u)\) ，换言之 \(s(u) + s^2(u) = u + s(u)\) ，或等价地

\[
s ^ {2} (u) = u\tag{11}
\]

于是 \(s^{2}\) 是 E 的恒等映射。由此可得

\[
\mathrm{T} (s (u)) = \mathrm{T} (u), \quad \mathrm{N} (s (u)) = \mathrm{N} (u).\tag{12}
\]

最后，关系式 \((u - u)(u - s(u)) = 0\) 给出

\[
u ^ {2} - \mathrm{T} (u). u + \mathrm{N} (u) = 0\tag{13}
\]

对所有 \(u \in \mathbf{E}\) 成立。

命题 4. 设 E 是一个 A-代数，s 与 \(s'\) 是 E 的反自同构，使得 \((E, s)\) 与 \((E, s')\) 是凯莱代数。若 E 有一个包含单位元 e 的基，则 \(s' = s\) 。

显然 \(s'(u) = s(u) = u\) 对所有 \(u \in Ae\) 成立。若 T, N（相应地，T', N'）是 (E, s)（相应地，(E, s')）的迹函数与范数函数，则由 (13) 可得

\[
(\mathrm{T} (u) - \mathrm{T} ^ {\prime} (u)). u - (\mathrm{N} (u) - \mathrm{N} ^ {\prime} (u)) = 0.
\]

设 B 是 E 的一个包含 e 的基，u 是 B 中异于 e 的一个元素；则 \(\mathrm{T}(u) - \mathrm{T}'(u) = 0\) ，从而 \(s(u) = s'(u)\) 。由于 s 与 \(s'\) 在 B 上重合，故二者相等。

在下文中我们将记 \(\bar{u} = s(u)\) ，于是

\[
\left\{ \begin{array}{l l} u + \bar {u} = \mathrm{T} (u), & u \bar {u} = \bar {u} u = \mathrm{N} (u), \quad \bar {\bar {u}} = u \\ \overline {{u + v}} = \bar {u} + \bar {v}, & \overline {{\alpha u}} = \alpha \bar {u}, \quad \overline {{u v}} = \bar {v}. \bar {u} \end{array} \right.\tag{14}
\]

其中 \(u, v\) 属于 \(\mathbf{E}, \alpha \in \mathbf{A}\) ；此外

\[
\mathbf {T} (e) = 2 e, \quad \mathbf {N} (e) = e.\tag{15}
\]

由公式

\[
\mathrm{T} (u v) = u v + \overline {{{{u v}}}} = u v + \bar {v}. \bar {u} = u v + (\mathrm{T} (v) - v) (\mathrm{T} (u) - u),
\]

我们推出

\[
u v + v u = \mathrm{T} (u) v + \mathrm{T} (v) u + (\mathrm{T} (u v) - \mathrm{T} (u) \mathrm{T} (v))\tag{16}
\]

由此，交换 \(u\) 与 \(v\) ，得

\[
\mathbf {T} (v u) = \mathbf {T} (u v).\tag{17}
\]

另一方面， \(\mathrm{N}(u+v)=(u+v)(\bar{u}+\bar{v})=\mathrm{N}(u)+\mathrm{N}(v)+\mathrm{T}(u\bar{v})\) ，从而

\[
\mathrm{T} (v \bar {u}) = \mathrm{T} (u \bar {v}) = \mathrm{N} (u + v) - \mathrm{N} (u) - \mathrm{N} (v).\tag{18}
\]

现在，把 (16) 中的 u 换成 \(\bar{u}\) 便得到

\[
\mathrm{T} (u \bar {v}) = \mathrm{T} (u) \mathrm{T} (v) + \bar {u} v + v \bar {u} - \mathrm{T} (u) v - \mathrm{T} (v) \bar {u} = \mathrm{T} (u) \mathrm{T} (v) - u v - \bar {v}. \bar {u};
\]

从而

\[
\mathrm{T} (v \bar {u}) = \mathrm{T} (u \bar {v}) = \mathrm{N} (u + v) - \mathrm{N} (u) - \mathrm{N} (v) = \mathrm{T} (u) \mathrm{T} (v) - \mathrm{T} (u v).\tag{19}
\]

最后，显然对所有 \(\alpha \in \mathbf{A}\) ，

\[
\mathrm{N} (\alpha u) = \alpha^ {2} \mathrm{N} (u);\tag{20}
\]

特别地 \(\mathbf{N}(2u) = 4\mathbf{N}(u)\) ，于是公式 (19) 给出

\[
(\mathrm{T} (u)) ^ {2} - \mathrm{T} (u ^ {2}) = 2 \mathrm{N} (u).\tag{21}
\]

显然 T 是 (Ae)-模 E 上的线性形式。由于 \((u, v) \mapsto \mathrm{T}(v\bar{u})\) 是该模上的双线性形式，由 (18) 与 (20) 可知 N 是一个二次型（参见 IX，§ 3，第 4 号）。

